\section{Discussion}
\label{sec:discussion}

% ── 4.1 Local-first design ───────────────────────────────────────────────────
\subsection{Local-first architecture for government and HPC environments}

The decision to build \textit{Query-GSAS} as a fully local system was driven
by the operational context of large-scale synchrotron and neutron
scattering user facilities.
At facilities such as the APS, NSLS-II, SNS, and their international
counterparts, beamline computers commonly operate within institutional
network security perimeters that prohibit or restrict outbound
connections to commercial cloud API endpoints. Other commercial and government facilities may limit internet access for security purposes. 
A local-inference architecture eliminates both concerns: all data
remain within the facility network regardless of the question asked.

The quantised Llama 3 8B inference model ($\approx$5\,GB) and
Qwen2.5-3B ($\approx$2\,GB) are pre-downloaded and cached;
subsequent queries require no internet connectivity.
The ingestion pipeline uses only publicly available documentation
pages and does not contact any commercial API.
This makes \textit{Query-GSAS} suitable for deployment on HPC login nodes
and beamline control workstations without requiring firewall
exceptions or API key management.

% ── 4.2 GSAS-II native integration ───────────────────────────────────────────
\subsection{GSAS-II integration}

\textit{Query-GSAS} has been successfully integrated into GSAS-II and is made available to users via a Help menu command. \textit{Circa} 200 lines of code were used to provide this integration, but almost all of that is used to confirm the presence of needed components and download them where needed. No additional changes to GSAS-II's core codebase were required
beyond the \texttt{import} statement and a Help menu binding. If desired, the full 
\texttt{gsas-query} package can also be installed separately and run alongside GSAS-II. 

% ── 4.3 Comparison with keyword search ────────────────────────────────────────
\subsection{Advantages over keyword search}

Existing GSAS-II documentation access relies primarily on
browser-based keyword search within the documentation pages.
Keyword search is effective when the user knows the exact term
used in the documentation (\emph{e.g.}\ ``Mustrain'' for microstrain
broadening) but fails when users describe a concept in different
language (\emph{e.g.}\ ``peak width parameters'' or ``strain
broadening'').

Dense retrieval with sentence embeddings addresses this vocabulary
mismatch by embedding both query and document in the same semantic
space.
Our evaluation (Table~\ref{tab:retrieval}) shows Recall@6 of 0.98,
meaning the relevant documentation section is returned in the top
six results for 39 of 40 test queries.
The remaining miss reflects a structural vocabulary collision: the
question ``read back Rwp and lattice parameters from a script'' matches
tutorial pages throughout the corpus because those quantities appear
in every Rietveld tutorial, overwhelming the scripting context.
Hybrid retrieval combining dense embeddings with BM25 keyword scoring
would be expected to improve Recall@1 further while maintaining
the local-deployment requirement.
The inline citation mechanism further reduces the cognitive load of
verifying answers: each factual claim is linked to its source,
allowing expert users to quickly confirm or question the synthesised
response.

% ── 4.4 Limitations ─────────────────────────────────────────────────────────
\subsection{Limitations and future directions}

\paragraph{Index currency.}
GSAS-II is under active development and its documentation is updated
frequently.
To access any newly documented information, the index database must be periodically rebuilt using \texttt{gsas-query --setup --reset} to incorporate new tutorials and
help pages or must be downloaded from the GitHub release. 
Likewise, when one of the \texttt{gsas-query} UI's opens a GSAS-II source documentation web page, the page that is opened will be the document that was identified at the time of indexing, and the contents may have moved to a different section, should the index be stale. 
Automatic generation of the index for GitHub download is performed weekly. We plan to investigate the possibility of automatic regeneration when the documentation is changed. 
When run from within the GSAS-II GUI, users are encouraged to update the index at least on a monthly basis. 
Because tutorial URLs are resolved dynamically from
\texttt{tutorialIndex.py} at ingest time, newly published tutorials
are picked up automatically in subsequent runs without modifying
\texttt{sources.py}. Likewise, textbook sections and chapters are located dynamically. Manual configuration changes to the sources.py file will be needed if additional pages/chapters are added to the GSAS-II Home page, Help pages or Programmers' manual.

\paragraph{PDF quality and future tooling.}
As noted previously, \texttt{gsas-query} can ingest PDF documents, but 
parsing quality is limited by the \texttt{pypdf} text extractor, which
struggles with multi-column layouts, mathematical notation, and
table structures. 
ML-based PDF-to-Markdown converters such as Marker
\citep{paruchuri2024marker} and Docling \citep{docling2024}
produce substantially cleaner output for scientific documents and
represent a natural upgrade path for the PDF ingestion pipeline.
Where a native HTML equivalent exists --- the HTML source is
always preferred.

\paragraph{Citation fidelity with small models.}
Open-weight models in the 3B--7B parameter range do not reliably follow
multi-citation system-prompt instructions.
The keyword-overlap citation injector provides a practical fallback
that achieves dense inline coverage without model-size constraints,
but it assigns citations based on lexical rather than semantic
matching, which may occasionally attribute a sentence to a
non-optimal source chunk.
Larger models (Llama 3 70B \citep{meta2024llama3}, Qwen2.5 72B \citep{qwen2025}) follow the citation
instruction reliably and do not require the injector; they are
accessible via Ollama on GPU-equipped workstations.

\paragraph{Multi-modal queries.}
GSAS-II tutorials contain numerous annotated screenshots and
figures that are not currently indexed.
Incorporating vision-language models to caption and embed figures
would improve answers to questions about specific dialog layouts
or GUI element locations.

\paragraph{Integration with APEXA.}
\textit{Query-GSAS} is designed for integration into APEXA
(Advanced Photon EXperiment Assistant), a broader agentic framework
for beamline data analysis developed at the APS.
Within APEXA, \textit{Query-GSAS} will be exposed as a tool-callable
Model Context Protocol (MCP) server, enabling the orchestrating
agent to retrieve GSAS-II documentation context automatically
when a refinement workflow step requires clarification.

\paragraph{Off-line Equation Rendering.}
As noted, the GUI version of \textit{Query-GSAS} will use the MathJax JavaScript library to render equations, but this library is accessed over the internet. With modifications to file \texttt{gui.py}, this library could be installed locally and then equations could be rendered even on an air-gapped computer. 